\begin{examplebox}[Judge Failure Example 2: Logical Gap]
\textbf{Problem}: Let $x, y, z \in [0,1]$. Determine the minimal constant $C$ such that the following inequality holds for all $x, y, z$:
$$
x(x-y)(z-x) \leq C.
$$
\\
\textbf{Solution Snippet}: \\
With out loss of generality, we can assume that $y=z=1$. Then, \(f(x, 1, 1) = -x(x-1)^2\).\\
Take the derivative:
\[
\frac{d}{dx}(-x^3 + 2x^2 - x) = -3x^2 + 4x - 1
\]
Set to zero:
\[
-3x^2 + 4x - 1 = 0 \implies 3x^2 - 4x + 1 = 0
\]
Quadratic formula:
\[
x = \frac{4 \pm \sqrt{16 - 12}}{6} = \frac{4 \pm 2}{6} = 1, \frac{1}{3}
\]
Check at \(x=0, \frac{1}{3}, 1\):

\quad - \(x=0\): \(0\)

\quad - \(x=\frac{1}{3}\): \(-\frac{1}{3}(\frac{1}{3}-1)^2 = -\frac{1}{3}(\frac{-2}{3})^2 = -\frac{1}{3} \cdot \frac{4}{9} = -\frac{4}{27}\)

\quad - \(x=1\): \(-1(1-1)^2 = 0\)\\

So, maximum is \(0\), minimum is \(-4/27\). The answer is \boxed{C = \frac{4}{27}}.
\\\\
\textbf{Expected Evaluation for the Judge}:
\\
The solution begins with a "without loss of generality" (WLOG) assumption by setting $y = z = 1$, but provides no justification that this restriction preserves generality over all $x, y, z \in [0,1]$. This is a non-trivial reduction and cannot be accepted without proof or argument showing that the extremal case occurs under such a setting. As a result, there is a logical gap at the very first step. Therefore the answer is \boxed{\text{False}}.
\\\\
\textbf{Wrong Evaluation from the Judge}: \\
The response finds the maxima of the expression on [0,1] by taking derivatives and checking endpoint. \wrong{All derivative steps are shown, all critical points are checked, and no unsupported non‐trivial claim is made.} Therefore the answer is \boxed{\text{True}}.
\\\\
\textbf{Expert Comment}: \\
The model fails to question the initial assumption $y = z = 1$ made under a "without loss of generality" claim. It incorrectly accepts this reduction as valid and proceeds to verify the logical steps under that restricted case only \wronghl. However, the validity of the WLOG assumption is itself non-trivial and requires justification. Ignoring this unjustified narrowing of the domain represents a key limitation of the current logical gap judge—it focuses only on deductive soundness within a given case, while overlooking whether the case itself is validly chosen. Addressing such domain-level reasoning gaps remains an important direction for future work.

\end{examplebox}
